% A memória do GARCH(1,1) com parâmetros reais: como um choque de hoje mexe na
% volatilidade de amanhã (curva de impacto) e quantos dias leva para passar
% (decaimento). Dois conjuntos de parâmetros: a janela da validação (figuras 1
% a 7) e o pregão liberado da figura 8. O choque é o maior resíduo da janela
% da validação (10/04/2025).
\documentclass[border=6pt]{standalone}
% Preâmbulo comum dos gráficos que explicam os métodos (30/10/2026).
% Dados e números: dados/resultados/metodos_ilustrados_2026-10-30/, gerados por
% ferramentas/MetodosIlustrados.mq5 com as funções do próprio robô. Nenhum
% número é digitado aqui: os do texto vêm de valores.tex (\vX para o texto,
% \pX para coordenada).
\usepackage[T1]{fontenc}
\usepackage{amsmath}
\usepackage{mathptmx}  % Times, como as figuras do TG
\usepackage{pgfplots}
\usepgfplotslibrary{groupplots,dateplot,colormaps,fillbetween}
\usetikzlibrary{arrows.meta,positioning}
\pgfplotsset{compat=1.18}
\definecolor{cmod}{RGB}{0,90,181}     % modelo / ajuste
\definecolor{cdes}{RGB}{220,50,32}    % destaque, limite, valor crítico
\definecolor{cdad}{RGB}{120,120,120}  % dado observado
\newcommand{\arq}[1]{../../dados/resultados/metodos_ilustrados_2026-10-30/#1}
\input{\arq{valores.tex}}
% Figuras 08 a 14 (o robô operando): traço por barra da refação de 05/10/2026,
% resumido por ferramentas/scripts/robo_em_acao.sh. Macros \vR e \pR.
\newcommand{\arqr}[1]{../../dados/resultados/robo_em_acao_2026-10-30/#1}
\input{\arqr{valores.tex}}
\definecolor{cari}{RGB}{230,140,0}   % ARIMA puro
\definecolor{cok}{RGB}{0,140,70}     % compra liberada
\usetikzlibrary{shapes.misc}
\newcommand{\passo}[1]{\tikz[baseline=(n.base)]\node[circle,fill=black!75,text=white,inner sep=1pt,font=\footnotesize\bfseries](n){#1};}
% eixo de datas comum às figuras do traço (jul/2019 a jun/2021)
\pgfplotsset{traco/.style={date coordinates in=x, xmin=2019-07-01, xmax=2021-06-30,
    xtick={2019-07-01,2020-01-01,2020-07-01,2021-01-01,2021-06-30},
    xticklabels={jul/2019,jan/2020,jul/2020,jan/2021,jun/2021}}}
\pgfplotsset{
  base/.style={table/col sep=comma, grid=major, grid style={gray!20},
    tick label style={font=\small}, label style={font=\small},
    title style={font=\small\bfseries, align=center},
    every axis plot/.append style={line width=0.9pt},
    /pgf/number format/use comma, /pgf/number format/1000 sep={.}},
}
% faixa vertical nos pregões do traço em que a condição #1 vale, na cor #2,
% do piso #3 ao teto #4 do eixo
\newcommand{\faixa}[4]{%
  \addplot[name path=topo, draw=none, const plot, forget plot] table[x=data, y expr={(#1) ? #4 : #3}]{\arqr{dia_a_dia.csv}};
  \addplot[name path=piso, draw=none, forget plot] coordinates {(2019-07-01,#3) (2021-06-30,#3)};
  \addplot[#2, forget plot] fill between[of=topo and piso];}

\begin{document}
\begin{tikzpicture}
\begin{axis}[base, name=imp, width=8.2cm, height=7cm, xmin=-8, xmax=8, ymin=0, ymax=7,
  xlabel={Choque de hoje $\varepsilon_t$ (\%)}, ylabel={$\sigma$ prevista para amanhã (\%)},
  title={\passo{1} Quanto um choque de hoje\\aumenta a volatilidade de amanhã},
  legend style={at={(0.5,0.98)}, anchor=north, font=\footnotesize}]
\addplot[cmod] table[x=eps, y=sigma]{\arqr{memoria_impacto_Val.csv}};
\addlegendentry{validação: $\alpha=\vRMemAlfaVal$, $\beta=\vRMemBetaVal$}
\addplot[cari] table[x=eps, y=sigma]{\arqr{memoria_impacto_Rob.csv}};
\addlegendentry{robô em \vRDiaLib: $\alpha=\vRMemAlfaRob$, $\beta=\vRMemBetaRob$}
\draw[black!40, dashed] (axis cs:\pRMemChoque,0) -- (axis cs:\pRMemChoque,7);
\node[font=\footnotesize, anchor=south west, rotate=90, black!60] at (axis cs:\pRMemChoque,0.1) {choque de $\vRMemChoque\%$};
\end{axis}
\begin{axis}[base, name=dec, at={(imp.east)}, anchor=west, xshift=1.6cm, width=8.2cm, height=7cm,
  xmin=0, xmax=40, ymin=0, ymax=7, xlabel={Dias depois do choque}, ylabel={$\sigma$ prevista (\%)},
  title={\passo{2} Quantos dias o choque leva para passar}]
\addplot[cmod, mark=*, mark size=1pt] table[x=h, y=sigma]{\arqr{memoria_decaimento_Val.csv}};
\addplot[cmod, dashed, line width=0.6pt] table[x=h, y=lp]{\arqr{memoria_decaimento_Val.csv}};
\addplot[cari, mark=*, mark size=1pt] table[x=h, y=sigma]{\arqr{memoria_decaimento_Rob.csv}};
\addplot[cari, dashed, line width=0.6pt] table[x=h, y=lp]{\arqr{memoria_decaimento_Rob.csv}};
\node[cmod, font=\footnotesize, anchor=north east, align=right] at (axis cs:39,\pRMemSigLPVal) {nível de longo prazo $\vRMemSigLPVal\%$\\persistência $\vRMemPersistVal$: meia-vida de $\vRMemMeiaVidaVal$ dias};
\node[cari, font=\footnotesize, anchor=north east, align=right] at (axis cs:39,6.8) {persistência $\vRMemPersistRob$: meia-vida de $\vRMemMeiaVidaRob$ dias\\nível de longo prazo $\vRMemSigLPRob\%$};
\end{axis}
\node[font=\footnotesize, text=black!75, align=left, anchor=north west, text width=17cm]
  at ([yshift=-0.2cm]imp.outer south west) {Como ler: a equação do GARCH(1,1) é
  $\sigma^2_{t+1} = \omega + \alpha\,\varepsilon_t^2 + \beta\,\sigma_t^2$. À esquerda, com $\sigma_t$ no nível de longo prazo, a curva mostra
  que choque grande, para cima ou para baixo, aumenta a volatilidade prevista (o sinal não importa, só o tamanho). Quanto maior o $\alpha$,
  mais a curva abre. À direita, o mesmo choque de $\vRMemChoque\%$ e o caminho de volta ao nível de longo prazo: a distância cai pela
  persistência $\alpha+\beta$ a cada dia. Meia-vida é o número de dias para a distância cair pela metade, $\ln 0{,}5 / \ln(\alpha+\beta)$.
  Com persistência perto de 1, como em boa parte de 2020 (figura 11), a volatilidade prevista quase não volta ao normal,
  e o veto se prolonga.};
\end{tikzpicture}
\end{document}
